% TLDR: Preserve exact historical scalar evidence and its exposed-test limitations.
\section{Historical Judge Evidence}
\label{sec:legacy}
% Provenance: experiments/archive/01_veribench_25_tasks_3_samples_5_raters_judge_evaluation_start_06Nov2026_end_06Nov2026/SOURCE_MANIFEST.json; source commit be287d9bb2f7bd35ba6db1b8735c9ba10d00eade.
The historical archive contains several distinct analyses that must not be pooled or described as one original result.
\subsection{Original audit and score mapping}
% Provenance: original VeriBench appendix G, Table 9; paper_original in SOURCE_MANIFEST.json; public NeurIPS_2026_VeriBench.pdf.
The original VeriBench report describes a raw Codex command-line rubric audit on 75 artifacts~\citep{veribench2026}.
It reports Pearson correlation $0.6137$, Spearman correlation $0.3755$ and Kendall correlation $0.2996$ against its historical pooled human target.
Two-fold isotonic score mapping in that audit gives $0.2108$, $0.0461$ and $0.0416$, respectively.
The source provides no confidence intervals for these entries.
That target pools seven rating files, including two earlier-prompt passes from one person, and differs from the later five-rater snapshot.
This full-file rubric audit is separate from VeriBench's deployed coverage scorer; its agreement cannot be transferred to a different judge implementation.

An early analysis reports Pearson correlation $0.703349$ on 72 scored artifacts from a 75-artifact, 25-task pool.
It uses a GPT-5.3-Codex judge with an isotonic score mapping fitted on training folds, against the mean of one expert's two earlier-prompt passes.
Its Spearman and Kendall values are $0.321143$ and $0.265060$.
These are two-fold task-split validation results without a separate test counterpart; the aggregate receipt does not supply confidence intervals.

\subsection{Selected scalar judge and agreement limits}
A later selected ensemble, source experiment 100, uses GPT-6 Sol with high reasoning effort and Claude Opus 5.5 with maximum compute effort through the historical Vals model gateway.
For each member it takes the median of three calls, then averages the two medians: six calls per completed scalar score.
It uses human rating instructions, training-derived conventions, and training examples.
The 25 tasks are split 10/5/10 into training, validation, and test, retaining each task's three candidates together.
The human target is the mean of available ratings from a three-expert panel, divided by five.
Of the 30 test artifacts, 18 have three expert ratings and 12 have two.

% Provenance: exp100_TEST_RECEIPT.json; all intervals below copied, not rerun.
\begin{table}[t]
\centering
\begin{tabular}{lrr}
\toprule
Statistic & Estimate & 95\% interval \\
\midrule
Spearman $\rho$ & 0.826 & [0.316, 0.966] \\
Kendall $\tau_b$ & 0.698 & [0.216, 0.913] \\
Pearson $r$ & 0.858 & [0.707, 0.941] \\
ICC(A,1) & 0.727 & [0.391, 0.890] \\
\bottomrule
\end{tabular}
\caption{Historical selected ensemble: 30 artifacts from ten test tasks.
Intervals are percentile task bootstraps with 10,000 requested draws.
ICC(A,1) is single-score absolute-agreement intraclass correlation.
The split had prior exposure; these are not fresh-study results.}
\label{tab:legacy}
\end{table}

The headline correlations in Table~\ref{tab:legacy} coexist with normalized mean absolute error $0.201111$ and mean bias $+0.181111$, with judge mean $0.653333$ and human mean $0.472222$.
Bias is the judge rating minus the expert mean, so a positive value indicates greater leniency.
% Provenance: POSTHOC_ERROR_INTERVALS_11-06-2026.json, newly computed from old data, not original receipt.
A separately labeled post hoc task bootstrap of the frozen historical data gives 95\% intervals [0.130, 0.273] for absolute error and [0.099, 0.262] for bias.
It uses 10,000 draws and reproduces the original point estimates; it is new analysis of old data, not new human evidence.
No new significance test is asserted for these quantities.

Excluding two test tasks whose references were placeholders when rated leaves 24 artifacts from eight tasks.
On that ordinary-reference subset, Spearman correlation is $0.652$ with interval [$-0.166$, $0.916$], and Pearson correlation is $0.835$ with interval [0.454, 0.958].
The mean bias rises to $+0.205556$, from judge mean $0.566667$ and human mean $0.361111$.
The subset correlation intervals use 1,998 nondegenerate draws out of 2,000 requested task bootstraps.
Reference quality therefore changes both the interpretation and the uncertainty of the headline association.

The reported mean repeat standard deviation, $0.020488$, also needs its original unit.
It is computed from three per-pass averages, each combining one call from each model.
It is not the standard deviation of repeated executions of the complete six-call median-based ensemble.
The archive supplies no interval for that repeat statistic.

For context, the archived leave-one-expert-out comparison reports Spearman correlation $0.764$ over all 75 artifacts and $0.747$ over the 45 validation-plus-test artifacts of experiment 100.
These compare one expert with the other available experts and have no uncertainty intervals in the imported aggregate.
Their task sets and targets differ from the 30-artifact test comparison in Table~\ref{tab:legacy}; they cannot establish that the selected judge exceeds human agreement on a matched test set.

Several limitations prevent clean transfer claims.
Candidate-generation records trace 74 of 75 files to three single-call, zero-shot samples per task from Claude 3.7 Sonnet at temperature $0.7$, with one source receipt missing.
These are repeated samples of one older model, not three different coding agents.
The source receipt discloses prior exposure to the test ratings in earlier experiments and an initial descriptive audit; the exact total number of exposed judge versions and resulting selection inflation are not established by the inspected records.
Nine references changed statement content between rating and earlier judging, although experiment 100 restored rating-era inputs.
The reconstruction remains ambiguous for five tasks whose references changed during the rating window; six test artifacts from two tasks carry that ambiguity flag.
The source appendix also reports three test items whose rating sheet linked a different candidate artifact.
Excluding either group raises test agreement in the source analysis, but does not repair the original evidence or justify replacing its full denominator.
Neither adding raters to those files nor narrowing their intervals would erase these provenance limits.

Finally, a separate five-fold automated recipe, source experiment 104, reports tie-preserving Spearman correlation $0.603$ [0.287, 0.836] over all 75 artifacts and 25 tasks, with Kendall $0.498$ and Pearson $0.621$.
Its original floating-point calculation gave Spearman $0.601533$; the published addendum preserves mathematical ties before ranking.
This evaluates fold-specific recipe fits, not the frozen selected ensemble.
The archive retains both versions and their labels.
These original analyses concern scalar LLM judges; they do not evaluate the proposed native probability interface or validate a prover judge.

\subsection{Retrospective comparison of training objectives}
% Provenance: experiments/archive/10_fit_5_judge_mixtures_to_75_outputs_pearson_vs_ranked_probability_start_09Nov2026_end_09Nov2026/results.md and results/results.json; existing analysis only.
A subsequent matched comparison fits mixtures of five cached judge configurations on the same historical ratings.
Three scalar repeats per configuration and artifact are converted to floored six-category histograms; no new semantic-model calls are made for this comparison.
Both training objectives use the same mixture class and select shrinkage by validation RPS.
Its original task split yields 30 training, 15 validation and 30 test artifacts, with 78, 39 and 78 observed expert ratings.

On the 30 test artifacts from ten tasks and nine task families, the normalized RPS difference, RPS-trained minus Pearson-trained, is $+0.005241$.
The 95\% family-bootstrap interval is [$-0.068074$, $+0.014864$], using 1,998 valid refit/reselection draws out of 2,000 requested.
Lower RPS is better; the point estimate favors Pearson, but the interval does not resolve an advantage or establish equivalence.
Validation selected the RPS-trained mixture, and this selection was not changed after test inspection.
No significance test is asserted.

The bootstrap refits and reselects mixtures but conditions on the cached model calls, prompts and rater identities.
It retains historical task-family overlap and prior exposure rather than repairing them.
Additional cached-metric, cross-validation and expanded-judge-pool analyses reuse this human evidence; a separate synthetic study contains no human observations.
They motivate the prospective comparison but do not constitute independent replications or fresh-rating validation.
